% ============================================================
%  8. Future Work
% ============================================================
\section{Future Work}
\label{sec:future}

\paragraph{CRF decoding.}
Replacing the softmax output layer with a Conditional Random
Field~\citep{lafferty2001crf} would enforce valid BIO transition constraints
globally.  This is particularly beneficial under high noise, where isolated
\texttt{I-X} predictions without a preceding \texttt{B-X} become more common.

\paragraph{Multilingual extension.}
\RetVec{} is language-agnostic by design: it hashes Unicode code points rather
than looking up a vocabulary.  Applying the same distillation pipeline with
mBERT~\citep{devlin2019bert} or XLM-R~\citep{conneau2020unsupervised} as
teacher would yield a single multilingual \LightRet{} without per-language
vocabularies, directly applicable to cross-lingual NER transfer.

\paragraph{Noise curriculum learning.}
The current Stage~3 noise schedule applies fixed perturbation probabilities
throughout training.  Progressive noise curricula~\citep{bengio2009curriculum}
---starting from clean text and gradually increasing perturbation intensity---may
improve convergence and final robustness, particularly in the high-noise regime.

\paragraph{CRF-augmented label projection.}
The dynamic BIO projection (Section~\ref{subsec:label_proj}) uses a
deterministic rule for continuation labels.  A learned projection that
scores all possible label sequences for a merged or split group jointly
would reduce projection errors in complex cases (e.g., space insertion
inside a multi-word named entity).

\paragraph{Structured prediction with noisy output.}
Stage~3 currently treats noisy-word NER as a flat sequence labeling task.
Span-based decoders~\citep{joshi2020spanbert,shen2021locate} could be
explored, as they are less sensitive to word boundary shifts and naturally
handle overlapping or discontiguous mentions.

\paragraph{On-device deployment.}
At ${\sim}4$M parameters, \LightRet{} occupies ${\sim}16$\,MB in float32.
Post-training quantisation to INT8 would reduce this to ${\sim}4$\,MB,
enabling deployment in mobile NLP pipelines and edge devices where BERT-scale
models are impractical.

\paragraph{Domain adaptation.}
The distillation corpus (WikiText-103 + CoNLL-2003) is news-domain heavy.
Domain-adaptive distillation~\citep{xu2020bert} on clinical, legal, or
social-media text is a natural next step, especially given \LightRet{}'s
vocabulary-free backbone which does not suffer from OOV degradation in
specialised terminologies.
